% Insert after the existing complete-exemplar/finite-visibility subsection.
\subsection{When a replay gradient improves a neural exemplar}
\label{app:theory-neural-replay-response}

Lemma~\ref{lem:shared-exemplar-retention} permits shared parameters by
assuming descent of a joint potential. A local statement can instead expose
the condition under which a realized replay update helps a particular
exemplar. This condition involves interference among scores and the optimizer's
response, quantities that the categorical abstraction suppresses.

\begin{proposition}[Local replay response under shared parameters]
\label{prop:neural-replay-local-response}
For a selected bank of $k$ complete exemplars, define
\[
 \ell_b(\theta)=\frac{\log\pi_\theta(e_b\mid x_b)}{L_b},\qquad
 h_b=\nabla_\theta\ell_b(\theta),\qquad
 \bar h=\frac1k\sum_{j=1}^k h_j.
\]
Suppose a step has the form
$\Delta=\eta H(g+\lambda\bar h)+u$, where $\eta,\lambda\ge0$,
$H$ is a fixed symmetric positive semidefinite matrix for this step,
$g$ is the fresh ascent direction, and $u$ is the remaining update.
Assume $\ell_b$ is twice differentiable and its Hessian has operator norm
at most $B_b$ along $\{\theta+s\Delta:0\le s\le1\}$.
Set
\[
 K_{bj}=h_b^\top Hh_j,\qquad
 \kappa_b=\frac1k\sum_jK_{bj},\qquad
 m_b=\eta h_b^\top Hg+\eta\lambda\kappa_b
       +h_b^\top u-\frac{B_b}{2}\|\Delta\|_2^2.
\]
Then
\begin{equation}
 \ell_b(\theta+\Delta)-\ell_b(\theta)\ge m_b,
 \qquad
 \pi_{\theta+\Delta}(e_b\mid x_b)
    \ge\pi_\theta(e_b\mid x_b)e^{L_bm_b}.
 \label{eq:neural-replay-local-response}
\end{equation}
In particular, $m_b>0$ guarantees that this update increases the complete
exemplar's probability, whether or not it appeared among the fresh samples.
\end{proposition}

\begin{proof}
Taylor's theorem with the Hessian bound gives
$\ell_b(\theta+\Delta)-\ell_b(\theta)
 \ge h_b^\top\Delta-(B_b/2)\|\Delta\|_2^2$.
Substitution of the stated step and
$h_b^\top H\bar h=k^{-1}\sum_jK_{bj}$ proves the first inequality.
Multiplication by $L_b$ and exponentiation prove the second.
\end{proof}

The required replay response is a row average $\kappa_b$ of the score Gram
matrix, not only its nonnegative diagonal entry $K_{bb}$. A sufficient
interference condition is $\kappa_b>0$, with replay strength satisfying
\[
 \eta\lambda\kappa_b>
   -\eta h_b^\top Hg-h_b^\top u
       +\frac{B_b}{2}\|\Delta\|_2^2.
\]
This inequality includes the step's curvature cost, which itself depends on
the chosen dose; arbitrarily increasing $\lambda$ is not a sufficient rule.

The interference condition cannot be dropped. For example, consider three
categorical logits sharing one scalar parameter, with slopes $(1,-2,0)$.
At a point with probabilities $(s,s,1-2s)$, where $0<s<1/4$, let the first
two categories be correct and banked. Their unit-length log-score gradients
are $h_1=1+s$ and $h_2=-2+s$. Uniform replay alone follows
$\bar h=s-1/2$, so
\[
 \left.\frac{d}{dt}\log p_1\right|_{\dot\theta=\bar h}
   =(1+s)(s-1/2)<0.
\]
Thus shared parameters can make replay initially reduce the probability of
a rare stored mode while improving the average replay score.

For an AdamW step, replay can alter both moment estimates. The proposition
therefore does not identify its $H$ with an unchanged AdamW preconditioner.
It applies to a specified decomposition of the realized step, with any
remaining contribution included in $u$; alternatively its Taylor bound can
be evaluated directly from $h_b^\top\Delta$. Neither the interference margin
nor a curvature bound is established for the experimental runs here.
The statement concerns complete stored responses under the sampling law
specified in App.~\ref{app:theory-exemplar-bridge}. Raising one such
probability strengthens that exemplar's contribution to its key, but need
not increase total key probability or transfer to a held-out prompt.
